%!TEX root = ../dissertation.tex

% =====================================================================
%  APPENDIX C — Online inference-time measurement
% ---------------------------------------------------------------------
%  Numbers from scripts/benchmark_inference_time.py in the thesis-cir repo
%  (100 queries sampled seed 42 from the instance_sample20_seed42 subset,
%  K=5 captions INSTR_F..J, max_new_tokens=256, 12 tiles, bf16).
%  Hardware: single NVIDIA RTX 6000 Ada (48 GB). Every stage timed with
%  torch.cuda.synchronize(); 3 warmup queries discarded. Results:
%  runs/inference_time/{per_query,summary}.csv.
% =====================================================================

\chapter{Online Inference-Time Measurement}\label{app:inference_time}

This appendix reports the \emph{online} cost of answering a single query with
the \emph{Ours} pipeline, and discusses where that cost can be reduced. The
database image embeddings are computed \emph{offline}, once, before any query is
issued, and are therefore excluded. For each query the timed work is: (i)
\textbf{caption generation}---the InternVL3.5-8B captioner runs the $K{=}5$
instruction series (\textsc{Instr\,F}--\textsc{J}) on the query image and
modification text---followed by (ii) \textbf{triplet embedding}---encoding the
query image, the modification text, and the $K{=}5$ generated captions with the
retrieval backbone. We report both backbones studied in the main text, CLIP
ViT-L/14 and SigLIP2. Caption generation is backbone-independent, so it is
shared between the two totals.

We time $100$ queries sampled from the \texttt{instance\_sample20\_seed42}
subset on a single NVIDIA RTX 6000 Ada GPU, wrapping every stage in a CUDA
synchronisation and discarding three warm-up queries. Table~\ref{tab:inference_time}
summarises the per-query distribution.

\begin{table}[h!]
\centering
\sisetup{table-format=5.1, group-minimum-digits=4}
\begin{tabular}{l S S S S S S}
\toprule
Stage (\si{\milli\second}) & {mean} & {min} & {Q1} & {median} & {Q3} & {max} \\
\midrule
Caption generation \emph{(shared)} & 14702.6 & 10138.1 & 12172.5 & 14718.4 & 16664.4 & 21341.7 \\
\midrule
CLIP-L: triplet embedding          &    42.7 &    17.8 &    19.7 &    21.5 &    26.0 &   257.9 \\
CLIP-L: \textbf{total}             & 14745.3 & 10156.0 & 12193.3 & 14749.6 & 16750.4 & 21363.8 \\
\midrule
SigLIP2: triplet embedding         &    67.3 &    42.6 &    44.3 &    45.5 &    50.0 &   281.2 \\
SigLIP2: \textbf{total}            & 14769.9 & 10180.9 & 12216.8 & 14774.8 & 16777.4 & 21388.1 \\
\bottomrule
\end{tabular}
\caption[Per-query online inference time (\emph{Ours}, $K{=}5$).]{Per-query
online inference time in milliseconds ($n{=}100$ queries, RTX 6000 Ada). The
database embedding is precomputed offline and excluded. Caption generation is
identical for both backbones.}
\label{tab:inference_time}
\end{table}

Two observations follow. First, \textbf{caption generation dominates}: it
accounts for roughly $99.7\%$ of the per-query cost ($\approx\!14.7$\,s median),
while embedding the triplet takes only tens of milliseconds. Second, the
\textbf{retrieval backbone is essentially irrelevant to latency}: SigLIP2 embeds
about twice as slowly as CLIP-L ($45.5$ vs.\ $21.5$\,\si{\milli\second} median),
but this $\approx\!24$\,\si{\milli\second} difference is negligible against a
$\approx\!14.7$\,s query, so the two totals are practically indistinguishable.
The choice of backbone can therefore be made purely on retrieval quality without
a latency penalty.

\section*{Optimisation Opportunities}

Because the captioner is the sole bottleneck, all meaningful speed-ups target it;
we did not pursue them here since latency is not a contribution of this thesis,
but several are readily available:

\begin{itemize}
  \item \textbf{Fewer captions.} We generate $K{=}5$ captions per query and
    average them (\texttt{avg-5}). Dropping to a single caption cuts the dominant
    cost by roughly $5\times$; the single-instruction variant
    (Table~\ref{tab:res_ladder}) is only modestly behind \texttt{avg-5}, so this
    trades a small quality loss for a large latency gain.
  \item \textbf{Shorter generations and fewer tiles.} Caption-generation time
    scales with the number of decoded tokens (\texttt{max\_new\_tokens}$=256$)
    and the number of image tiles (up to $12$); both can be reduced with limited
    impact on caption content.
  \item \textbf{A lighter or faster captioner.} A smaller or distilled VLM,
    weight quantisation (INT8/INT4), a fused attention kernel (FlashAttention,
    disabled in our runs), or speculative decoding would each lower the
    per-token cost without changing the pipeline.
  \item \textbf{Batching.} The five instructions are currently issued
    sequentially; batching them (and batching across concurrent queries) would
    amortise fixed per-call overhead on modern GPUs.
\end{itemize}
